\documentclass[12pt,a4paper]{article}
\usepackage[utf8]{inputenc}
\usepackage[T1]{fontenc}
\usepackage{enumitem}
\usepackage[margin=2.5cm]{geometry}
\usepackage{parskip}
\usepackage{amsmath}
\usepackage{xcolor}

\title{final try — Final Exam --- Answer Key}
\date{}

\begin{document}

\maketitle

{\large \textbf{\textcolor{red}{ANSWER KEY --- Not for Distribution}}}

\vspace{0.5em}



\noindent
\textbf{Total points:} 6
\hfill
\textbf{Number of questions:} 6

\vspace{1em}
\hrule
\vspace{1.5em}


\noindent
\textbf{Question 1 (2.00 pts).}~
Which of the following is equivalent to a neural network with a nonlinear output?


\begin{enumerate}[label=\textbf{\Alph*.}]

  \item A linear combination of inputs

  \item \textbf{\textcolor{teal}{[CORRECT]}}  A weighted sum of inputs passed through an activation function

  \item A simple threshold value

  \item A random combination of inputs

\end{enumerate}



\textit{\small Explanation: A neural network is defined as a nonlinear model, where the output is a nonlinear combination of inputs.}



\vspace{0.8em}
\noindent\rule{\linewidth}{0.2pt}
\vspace{0.3em}


\noindent
\textbf{Question 2 (1.00 pts).}~
Neural networks are able to generalise.


\textbf{Answer:} \textcolor{teal}{True}



\textit{\small Explanation: This is supported by the fact that neural networks are able to learn from a large number of parameters and make predictions on new, unseen data.}



\vspace{0.8em}
\noindent\rule{\linewidth}{0.2pt}
\vspace{0.3em}


\noindent
\textbf{Question 3 (1.00 pts).}~
How would you represent the output for an image of the digit 5 in a classification problem with 10 classes, one for each digit from 0 to 9?


\textbf{Model Answer:} \textcolor{teal}{The output would be a vector with 10 elements, where the 6th element is 1 and all other elements are 0, i.e., (0,0,0,0,0,1,0,0,0,0).}



\textit{\small Explanation: 10 classes, one for each digit; One-hot encoding for the output}



\vspace{0.8em}
\noindent\rule{\linewidth}{0.2pt}
\vspace{0.3em}


\noindent
\textbf{Question 4 (1.00 pts).}~
How do the advancements in computer technology from the third generation to the fourth generation impact the user experience and the functionality of computers?


\textbf{Model Answer:} \textcolor{teal}{The advancements in computer technology from the third generation to the fourth generation allow users to interact with computers through keyboards and monitors, and enable the device to run many different applications at one time with a central program that monitors the memory. Additionally, the development of the microprocessor leads to the creation of small, powerful computers that can be networked, eventually leading to the development of the Internet.}



\textit{\small Explanation: Advancements in computer technology improve user experience; Fourth generation computers enable multitasking and networking}



\vspace{0.8em}
\noindent\rule{\linewidth}{0.2pt}
\vspace{0.3em}


\noindent
\textbf{Question 5 (1.00 pts).}~
Using the provided relationship curves between TP and Temp at a variety of constant CHL, and the relationship curves between TP, Secchi, and TSS at a constant CHL, evaluate how changes in Total Phosphorus (TP) might impact water transparency or clarity as measured by Secchi disk depth, considering the role of Total Suspended Solids (TSS) and Chlorophyll (CHL).


\textbf{Key Points:} \textcolor{teal}{I. Introduction to the relationship between TP, Secchi, and TSS. II. Analysis of the impact of TP on water transparency using the provided curves. III. Discussion of the role of TSS and CHL in this relationship. IV. Conclusion on the implications for water quality.}



\textit{\small Explanation: Consider the impact of TP on water quality and eutrophication, and how this relates to the provided relationship curves.}



\vspace{0.8em}
\noindent\rule{\linewidth}{0.2pt}
\vspace{0.3em}


\noindent
\textbf{Question 6 (1.00 pts).}~
Compare and contrast the characteristics of supervised and unsupervised learning in data mining, using specific examples from the provided context to support your analysis.


\textbf{Key Points:} \textcolor{teal}{I. Introduction to supervised and unsupervised learning in data mining
II. Characteristics of supervised learning: desired output, input, and output types
III. Characteristics of unsupervised learning: desired output, input, and output types
IV. Comparison of supervised and unsupervised learning: types of problems suited for each approach
V. Case study: Cluster Example and Compression Example
VI. Conclusion}



\textit{\small Explanation: Address the differences in desired output, input, and output types, as well as the types of problems each approach is suited for, using the Cluster Example and Compression Example as case studies.}





\end{document}
